\documentclass{situationist-preprint}

\title{The Upload Is a Successor}
\msid{PID-14-a}
\msfamily{PID-14}
\msdiscipline{Personal Identity and the Self}
\mstype{Principal manuscript}
\mskeywords{mind uploading; whole brain emulation; personal identity; gradual replacement; successor; pattern identity}

\begin{document}
\maketitle

\begin{abstract}
Would uploading your mind to a computer preserve you? We argue that it would not: uploading creates a psychologically continuous successor, not the original person. Two arguments support the conclusion. The first is from PID-05: you are a thinking organism, and the computer that runs the upload is not that organism. The second does not depend on animalism. Psychological continuity sustains the persistence of a person only when its later stages are realized in a thing of the same kind as its earlier ones, so that the stages are linked by a single persisting dynamics, as the composition principle of MET-06-a requires of the temporal parts of any persisting object. Teletransportation (PID-02) meets this condition; uploading does not. Nondestructive uploading fails in addition by branching. The upload is nonetheless a person if thinking is functional (MND-08), and it has much of what matters in survival (PID-06). Gradual neural replacement that keeps the organism alive is not uploading, and may preserve you.
\end{abstract}

\section{The problem}

The proposal that a human mind might be transferred to a computer was given currency by \citet{Moravec1988} and \citet{Kurzweil2005}, and a research programme for whole brain emulation was set out by \citet{SandbergBostrom2008}. Uploading is usually imagined in three forms. In \emph{destructive} uploading, the brain is scanned in detail, destroyed in the process, and an emulation is run on a computer. In \emph{nondestructive} uploading the brain survives the scan. In \emph{gradual} uploading, neurons are replaced one by one by functionally equivalent devices, or their functions are transferred piecemeal to an external system.

\citet{Chalmers2010} argued that gradual uploading is the most promising route to survival, since there seems to be no point in a gradual process at which the person ceases, and was more cautious about destructive and nondestructive uploading. \citet{CorabiSchneider2012} and \citet{Schneider2019} argued that the upload would be at best a copy. The question has two parts: whether the upload would be a person, and whether it would be you. We answer yes to the first, under a condition, and no to the second.

\section{Framework}\label{s:frame}

\begin{definition}\label{d:upload}
An \emph{upload} of a person is a computational system that, after a scan of the person's brain, realizes a psychology functionally equivalent to the person's at the time of the scan and continuous with it.
\end{definition}

\begin{definition}\label{d:succ}
A \emph{successor} of a person is a later person psychologically continuous with her who is not she.
\end{definition}

\begin{principle}[Psychological criterion]\label{p:pc}
A later person is the same person as an earlier one if and only if they are psychologically continuous, by a non-branching process or, where the process branches, with bodily continuity as the tie-breaker (PID-01), through a cause of the right kind.
\end{principle}

\begin{principle}[Kind condition]\label{p:kind}
Psychological continuity is through a cause of the right kind only if the later stages are realized in a thing of the same kind as the earlier ones: for a human person, a living human organism.
\end{principle}

\begin{assumption}\label{a:func}
An upload is psychologically continuous with the original and, if thinking is a functional property (MND-08), is a person.
\end{assumption}

\begin{assumption}\label{a:notorg}
The computational system that realizes an upload is not a living human organism.
\end{assumption}

Principle~\ref{p:kind} is the substantive premise, and it needs a defense, since \citet{Parfit1984} distinguished a narrow version of the psychological criterion, requiring the normal cause, from a wide version allowing any reliable cause, and the kind condition lies between them. It is supported in three ways. First, it applies to persons the general condition on persistence that MET-06-a states for composition over time and that PID-11 adopts for persons: only temporal parts linked by the relevant dynamics compose a persisting object, in that later stages are what earlier stages evolve into under the laws governing them. A human person's psychology is realized in, and evolves under, the dynamics of a living brain in a living body. The stages of an upload evolve under a different dynamics, the execution of a program on hardware; they are produced from the earlier stages, but they do not continue them under the same laws. Second, the condition explains the verdict of PID-02-a, which a wide criterion also explains, while distinguishing teletransportation from uploading, which a wide criterion cannot: the teletransporter's replica is a living human organism, whose psychology evolves under the same dynamics as the original's, and the gap in transmission is bridged by a process that reconstructs a thing of the original's kind. Third, it is the condition under which REL-08-a observes that a post-mortem realizer could be the person on the pattern theory: a living organic body of the same kind.

\section{The argument}\label{s:arg}

\begin{theorem}\label{t:notperson}
No upload is the original person.
\end{theorem}

\begin{proof}
By Assumption~\ref{a:notorg} the upload is realized in a thing that is not a living human organism. By Principle~\ref{p:kind}, the psychological continuity between the original and the upload is therefore not through a cause of the right kind. By Principle~\ref{p:pc}, the upload is not the original person.
\end{proof}

\begin{theorem}\label{t:succ}
An upload is a psychologically continuous successor of the original, provided thinking is functional.
\end{theorem}

\begin{proof}
By Assumption~\ref{a:func} the upload is a person psychologically continuous with the original, and by Theorem~\ref{t:notperson} it is not the original. By Definition~\ref{d:succ} it is a successor.
\end{proof}

\begin{proposition}\label{p:animal}
If PID-05 is true, the upload is not you, independently of Principle~\ref{p:kind}.
\end{proposition}

\begin{proof}
By PID-05 you are a thinking organism, which persists only as long as its life continues. The upload is realized in a system that is not a living organism (Assumption~\ref{a:notorg}), and so is not the organism; in destructive uploading the organism dies or is left without a brain, and in nondestructive uploading it remains where it was.
\end{proof}

\begin{proposition}\label{p:branch}
In nondestructive uploading, the upload is not the original person even on a wide criterion that omits Principle~\ref{p:kind}.
\end{proposition}

\begin{proof}
If the brain survives the scan, the process branches: the original body and brain, and the upload, are both psychologically continuous with the person at the time of the scan. By Principle~\ref{p:pc} bodily continuity breaks the tie, and it favors the original. This is the structure of the Branch-Line case of PID-02-a.
\end{proof}

The two arguments, from the kind condition and from animalism, are independent and agree. The first concerns the person in the sense of PID-01; the second concerns the organism that, according to PID-05-a, each of us fundamentally is. They come apart in teletransportation, in which the person, but not the organism, survives; they agree that the upload is neither.

\section{Gradual replacement}\label{s:gradual}

\citet{Chalmers2010} considers a process in which a person's neurons are replaced, one at a time, by silicon devices that duplicate their input--output function, with the person conscious throughout. If at no step does the person cease, and the end state is a fully artificial brain, uploading by degrees seems to preserve the person.

We distinguish two cases. In the first, the devices are implanted in the living brain and integrated into the organism, which continues its life: they are supplied by its circulation, regulated by its physiology and share its fate. Then by MET-06-a they become parts of the organism, as prosthetic joints and pacemakers do, and Principle~\ref{p:kind} is satisfied throughout, since the psychology remains realized in a living human organism. The person survives, and on PID-05 so do you. This is not uploading in the sense of Definition~\ref{d:upload}, but prosthesis. In the second case, the functions are transferred step by step to an external computational system, and the biological brain is dismantled. At each step part of the psychology is realized outside the organism; at the end all of it is. The process is gradual, but its end point is an upload, and by Theorem~\ref{t:notperson} the end point is not the original. That there is no sharp step at which the original ceases shows that the transition is vague, not that the original survives; a person can cease gradually, as she can come into existence gradually.

\section{Objections and replies}\label{s:obj}

\begin{objection}[Parity with teletransportation]
If you survive being scanned, destroyed and rebuilt from new matter on Mars (PID-02-a), you survive being scanned, destroyed and run on new hardware. The difference between organic and silicon realizers is a difference of substrate, and substrate does not matter if thinking is functional.
\end{objection}

\begin{reply}
That thinking is functional (MND-08) settles whether the upload thinks, not whether it is you. Persistence conditions are not fixed by what a thing does at a time but by what makes stages at different times stages of one thing, and Principle~\ref{p:kind} identifies this with continuity under a single dynamics. The replica on Mars is a human organism whose psychology continues under the dynamics of the original's; the upload's does not. The difference is not one of substrate alone; it is a difference in the kind of process in which the psychology is realized.
\end{reply}

\begin{objection}[Pattern identity]
A person is a pattern of information, and wherever the pattern is realized, there is the person \citep{Kurzweil2005}.
\end{objection}

\begin{reply}
REL-08-a shows that a pattern theory of persons is coherent, and that if PID-05 is true it is not the true theory of us: a realizer of your pattern would have much of what matters but would not be you. The present argument adds that, even for persons as such, an unrestricted pattern theory conflicts with the composition principle of MET-06-a, which denies that arbitrary sums of stages compose persisting things. A pattern theory restricted to realizers of the original's kind is compatible with Theorem~\ref{t:notperson}.
\end{reply}

\begin{objection}[Then uploading is death]
If the upload is not you, destructive uploading is as bad as death, and no one should accept it.
\end{objection}

\begin{reply}
The inference assumes that identity is what matters. PID-03 and PID-06 argue that it is not, and that psychological connectedness carries most of what matters. An upload is strongly connected with the original, and so has much of what matters in her survival. Destructive uploading is therefore not as bad as death, though it is the end of the original person and of the organism she is. PID-14-b examines what the original has reason to want for her successor.
\end{reply}

\begin{objection}[Arbitrary kinds]
The kind condition is arbitrary at the margins: a person whose brain is half prosthetic, or who lives in a body sustained by machines, is of no clear kind.
\end{objection}

\begin{reply}
The condition is vague, as is the condition of being a living organism, and there will be cases in which it is unclear whether it holds. But a vague condition is not an arbitrary one, and uploading, which moves the psychology entirely into a computational system, is not a borderline case.
\end{reply}

\section{What the result does not establish}\label{s:not}

The result does not establish that uploads are persons; Theorem~\ref{t:succ} is conditional on MND-08, and ETH-24 holds that current artificial systems lack the capacities on which moral status depends. It does not establish that uploading is technically possible. It does not show that disembodied persons are impossible, which PID-10 denies, only that a human person does not become one by uploading. And it does not say that the upload's life would be worse than the original's.

\section{Conclusion}

Uploading would not preserve you. The upload, if thinking is functional, is a person psychologically continuous with you; but its psychology is realized in a thing of another kind, under another dynamics, so that the continuity is not of the kind that makes one person; and if animalism is true, it is not the organism you are. It is a successor, with much of what matters in your survival and without your identity.

\begin{thebibliography}{99}\small
\bibitem[Chalmers(2010)]{Chalmers2010} Chalmers, D.~J. (2010). The singularity: A philosophical analysis. \emph{Journal of Consciousness Studies}, 17(9--10), 7--65.
\bibitem[Corabi and Schneider(2012)]{CorabiSchneider2012} Corabi, J., and Schneider, S. (2012). The metaphysics of uploading. \emph{Journal of Consciousness Studies}, 19(7--8), 26--44.
\bibitem[Kurzweil(2005)]{Kurzweil2005} Kurzweil, R. (2005). \emph{The Singularity Is Near: When Humans Transcend Biology}. Viking.
\bibitem[Moravec(1988)]{Moravec1988} Moravec, H. (1988). \emph{Mind Children: The Future of Robot and Human Intelligence}. Harvard University Press.
\bibitem[Parfit(1984)]{Parfit1984} Parfit, D. (1984). \emph{Reasons and Persons}. Oxford University Press.
\bibitem[Sandberg and Bostrom(2008)]{SandbergBostrom2008} Sandberg, A., and Bostrom, N. (2008). \emph{Whole Brain Emulation: A Roadmap}. Technical Report 2008-3, Future of Humanity Institute, Oxford University.
\bibitem[Schneider(2019)]{Schneider2019} Schneider, S. (2019). \emph{Artificial You: AI and the Future of Your Mind}. Princeton University Press.
\end{thebibliography}

\end{document}
